

%<*preambleTag>

\newcounter{savepage}


\tolerance=1
\emergencystretch=\maxdimen
\hyphenpenalty=10000
\hbadness=10000
\hyphenchar\font=-1
\sloppy

\pagenumbering{gobble}

\setlist{topsep=0cm, noitemsep}

\renewcommand{\chapterheadstart}{\vspace*{\beforechapskip}}

%%%%%%%%%%%%%%%%%%%%%%%%%%%%%%%%%%%%%%%%%%%%%%%%%%%%%%%%%%%%%%%%
%%%%%%%%%Espaçamento antes e depois de cada capítulo%%%%%%%%%%%%
%%%%%%%%%%%%%%%%%%%%%%%%%%%%%%%%%%%%%%%%%%%%%%%%%%%%%%%%%%%%%%%%

\setlength\beforechapskip{0pt}
\setlength\afterchapskip{30pt}
%---------------------------------------------------------------


\geometry{top=3cm, bottom=2cm, left=3cm, right=2cm}
\newcommand{\thickhat}[1]{\mathbf{\hat{\text{$#1$}}}}
\deffootnote[0.65em]{0.65em}{0em}{\textsuperscript{\makebox[0em][r]{\thefootnotemark\ }}}
\DeclareCaptionFont{mysize}{\fontsize{12}{0.0}\selectfont}
\captionsetup{font=mysize}

%%%%%%%%%%%%%%%%%%%%%%%%%%%%%%%%%%%%%%%%%%%%%%%%%%%%%%%%%%%%%%%%
%%%%%%%%%%%%%%%%%%%%PREÂMBULO%%%%%%%%%%%%%%%%%%%%%%%%%%%%%%%%%%%
%%%%%%%%%%%%%%%%%%%%%%%%%%%%%%%%%%%%%%%%%%%%%%%%%%%%%%%%%%%%%%%%
\frenchspacing 
\renewcommand{\familydefault}{\sfdefault}
\usepackage[alf, abnt-etal-list=3, abnt-emphasize=bf,abnt-last-names=bibtex, abnt-etal-cite=3]{abntex2cite}	% Citações padrão ABNT
\bibliographystyle{abntex2-alf}
\renewcommand{\bibname}{Referências}

\makeatletter
\hypersetup{
	pdftitle={\@title}, 
	pdfauthor={\@author},
	pdfsubject={\imprimirpreambulo},
	pdfcreator={LaTeX with abnTeX2}, 
	colorlinks=true,  
	linkcolor=black,
	citecolor=black, 
	filecolor=black, 
	urlcolor=black,
	bookmarksdepth=4
}

%TOC Depth
\hypersetup{bookmarksopen=true,bookmarksopenlevel=4}
\setcounter{tocdepth}{4}

\makeatother

\setlength{\parindent}{1.5cm} % Tamanho do parágrafo
\setlength{\ABNTEXcitacaorecuo}{4cm}
\setlength{\intextsep}{0cm} %Espaçamento entre texto e figura
\setlength{\parskip}{0cm} % Espaçamento entre um parágrafo e outro

\makeindex

\newcommand{\apreambulo}{Trabalho apresentado como requisito para obtenção do título de bacharel em~\acurso~do Centro Tecnológico de Joinville da Universidade Federal de Santa Catarina. \break \break \aorientador \break \break \acoorientador}
\newcommand{\acidade}{Joinville}

\newcommand{\auniversidade}{
	\begin{SingleSpace}
		UNIVERSIDADE FEDERAL DE SANTA CATARINA \\
		CENTRO TECNOLÓGICO DE JOINVILLE \\
		\MakeUppercase{CURSO DE~\acurso}\\
	\end{SingleSpace}
}
%---------------------------------------------------------------


%MakeNomenclature
\makenomenclature
\renewcommand*{\nomname}{LISTA DE SIGLAS}

%</preambleTag>
%<*capaTag>
\thispagestyle{empty}

\begin{capa}
	\begin{center}
		
		\auniversidade
		\vspace{4.5cm}
		 
		\MakeUppercase{\aaluno} \\
		\vspace{5.5cm} 
				
		\MakeUppercase{\atitulo} \\
		
		\vfill
		\acidade \\
		\aano
		
	\end{center}
\end{capa}


%SEGUNDA PAGINA
\thispagestyle{empty}
\begin{capa}
	\begin{center}
		
        
		\MakeUppercase{\aaluno}\\
		\vspace{11.60cm}
 
		\MakeUppercase{\atitulo}   
	\end{center}
	%\vspace{-0.75cm}
		\vspace{1.5cm}	
		 {		
			\SingleSpacing \hspace*{7.5cm} \begin{minipage}{8.5cm}
				{\apreambulo}
			\end{minipage}
		
		}
		\vfill
	\begin{center}
		\acidade \\
		\aano
	\end{center}
\end{capa}

\textual

%</capaTag>


%<*folhadeaprovacaoTag>
\thispagestyle{empty}
%TERCEIRA PAGINA
\begin{capa}
 
\begin{center}
			      	      
		\MakeUppercase{\aaluno}\\
		\vspace{1.0cm}
 
		\MakeUppercase{\atitulo}   
	\end{center}
	%\vspace{-0.75cm}
		\vspace{1cm}	
		 {		
			\SingleSpacing \hspace*{7.0cm} \begin{minipage}{8.5cm}
				{Este Trabalho de Conclusão de Curso foi julgado adequado para obtenção do título de bacharel em \acurso, na Universidade Federal de Santa Catarina, Centro Tecnológico de Joinville.}
			\end{minipage}
			\newline
            \SingleSpacing \hspace*{7.0cm} \begin{minipage}{8.5cm}
				{Joinville (SC), \diaAprovacao~ de \mesAprovacao~ de \aano.}
			\end{minipage}
        }
        \vspace{1cm}
        
    \newline
    \textbf{Banca Examinadora:}

    \begin{center}
	    \vspace{2cm}
	    
	    \begin{center}
            ________________________________
        \end{center}
        \vspace{-0.5cm}
        \begin{center}
            \aorientador
        \end{center}
        \vspace{-1cm}
        \begin{center}
            Orientador(a)
        \end{center}
        \vspace{-1cm}
        \begin{center}
            Presidente
        \end{center}

	\end{center}
	
	\begin{center}
        \vspace{1.5cm}
        
	    \begin{center}
            ________________________________
        \end{center}
        \vspace{-0.5cm}
        \begin{center}
            \membrobancaA
        \end{center}
        \vspace{-1cm}
        \begin{center}
            Membro(a)
        \end{center}
        \vspace{-1cm}
        \begin{center}
            \UnimembrobancaA
        \end{center}

	\end{center}
	
	\begin{center}
        \vspace{1.5cm}
        
	    \begin{center}
            ________________________________
        \end{center}
        \vspace{-0.5cm}
        \begin{center}
            \membrobancaB
        \end{center}
        \vspace{-1cm}
        \begin{center}
            Membro(a)
        \end{center}
        \vspace{-1cm}
        \begin{center}
            \UnimembrobancaB
        \end{center}

	\end{center}
\end{capa}
    
\textual
%</folhadeaprovacaoTag>


%<*listOfFiguresTag>
\listoffigures*
\newpage
%</listOfFiguresTag>

%<*listOfTablesTag>
\listoftables*
\newpage
%</listOfTablesTag>

%<*tableOfContentsTag>
\renewcommand*{\contentsname}{SUM\'{A}RIO}
\tableofcontents*
\newpage
%</tableOfContentsTag>


%<*nomenclatureTag>
\printnomenclature
\newpage
%</nomenclatureTag>

%<*quadroTag>
\pdfbookmark[0]{\listofquadrosname}{loq}
\listofquadros*
\cleardoublepage
%</quadroTag>

%<*numberingTag>
\setcounter{savepage}{\arabic{page}}
\pagenumbering{arabic}
\setcounter{page}{\thesavepage+1}
%</numberingTag>

